\section{Input data for Initial Distribution of Test Particles}
\label{sec2.7}

\medskip

\textbf{General Remarks}

\medskip

The primary source (and: initial distribution of test particles,
in time dependent mode) is given as a function $Q(i, \vr,t
, \vv)$. $ Q $ is the density of the probability
distribution from which the species index $ isp $, the starting
point $\vr $, the velocity vector $\vv $ and
the starting time $ t $ are sampled in subroutine LOCATE. $(isp,
\vr,t , \vv)$ is the state of the starting
particle, which then will be ``followed" (traced) in subroutine FOLNEUT or
FOLION. There are five primary types of spatiotemporal distributions for primary sources (``Strata") , namely
Point sources, Line sources (to be written), Surface sources,
Volume sources and ``Census sources". The first four are uniform in a time-interval (or time-independent),
and the fifth one is an initial distribution (in volume) at a given point in time.

As the transport equation (within each iterative step, if any) is linear, various sources of test
particles can be treated subsequently and the responses can then
be linearly superposed.
EIRENE can handle up to NSTRA (see: PARMUSR, \cref{sec3.1})
different strata and it prints
output from
each single stratum as well as from the sum over strata. The source strength
is prescribed for each stratum separately (input card: FLUX(ISTRA), \dots).
Subdividing the total source into such strata can be useful
to increase
the efficiency of the code, (``stratified source sampling", see \cref{sec1.2.1b} or any
textbook on Monte Carlo integration), or if
the contribution of such sub-sources is of interest by itself.

In order to facilitate sampling each stratum can be subdivided
further into a number (NSRFSI) of ``sub-strata". Random sampling is
done by firstly sampling the sub-stratum ISRFSI, and then
(conditional) the initial state of the test flight within this
sub-stratum.

Random numbers from the multivariate distribution $Q$ are
generated by a sequence of uni-variate conditional
distributions (see \cref{eq1.15a}).
Source sampling from $Q$ in EIRENE is based on
the (formal) decomposition
\begin{align}\label{samp-Q}
Q(\vr,t,isp,\vv) &= Q_1(\vr,t) \times Q_2(isp|\vr,t)
\times Q_3(\vv|isp,\vr,t) \nonumber \\
&= \sum_i\!\!\int d\vv \,Q(\vr,t,i,\vv)
\times \int d\vv \,Q(isp,\vv|\vr,t)
\times Q(\vv |isp,\vr,t)
\end{align}
i.e.\ the source sampling routines in EIRENE
(samvol.f, samsrf.f, etc.)
start by sampling the spatio-temporal birth point
coordinates $\vr_0,t_0$
from $Q_1(\vr,t)$ (obtained by integrating $Q$ over velocity space $\vv$ and
summing over all species), then next sampling the species index $i$
from the conditional distribution $Q_2(i) $, which is conditional on $\vr_0,t_0$ and integrated over velocity space,
and finally then sampling the velocity $\vv_0$ from $Q_3(\vv)$ conditional on species $i_0$ and space-time $(\vr_0,t_0)$.


To facilitate random sampling from $Q_k$ for a particular stratum $k$, each stratum can be further subdivided into a sum of ``substrata":
\begin{equation}
Q_k= \sum_j w_{k,j} Q_{k,j}  \quad \texttt{with normalisation} \quad \sum_j  w_{k,j} = 1
\end{equation}
There are NSRFSI substrata, and the weights $w_{k,j}, \ j = 1, \dots, NSRFSI$, are used for first determining a particular
substratum $j_0$ by sampling, and then sampling from distribution $Q_{k,j_0}(\vr,\vv,i,t)$ of stratum $k$ as described above.

The total CPU resources (and storage for the census array, see
\cref{sec2.13}) are distributed over the strata. This
distribution is controlled by some of the input flags in this
input block, as described below. EIRENE stops working on a
particular stratum if either the assigned CPU-time for this
stratum has been reached (Message: ``No further CPU time for this
stratum"), or if all requested histories have been calculated
(Message: ``All histories for this stratum completed") or, in case of
time dependent problems, if the
storage on the census array of this stratum has been filled up
(Message: ``Census array filled for this stratum").

\medskip

\textbf{The Input Block}

\medskip

\begin{verbatim}
 *** 7. Data for primary sources
 NSTRAI
 INDSRC(ISTRA), ISTRA=1,NSTRAI
 ALLOC  AMPTS
 DO 70, ISTRAI=1,NSTRAI
   TXTSOU
   NLAVRP  NLAVRT  NLSYMP  NLSYMT
   NPTS  NINITL  NEMODS  NAMODS  NMINPTS
   FLUX  SCALV   IVLSF  ISCLS  ISCLT ...
   ...continue...ISCL1  ISCL2  ISCL3  ISCLB  ISCLA
 *Species index distribution
   NLATM  NLMOL  NLION  NLPLS NLPHOT
   NSPEZ
 *Distribution in physical space and time
   NLPNT  NLLNE  NLSRF  NLVOL  NLCNS
   NSRFSI
     DO 75, J=1,NSRFSI
       INUM    INDIM   INSOR   INGRDA(1)  INGRDE(1)
                               INGRDA(2)  INGRDE(2)
                               INGRDA(3)  INGRDE(3)
       SORWGT  SORLIM  SORIND  SOREXP  SORIFL
       NRSOR   NPSOR   NTSOR   NBSOR   NASOR   NISOR
       SORAD1  SORAD2  SORAD3  SORAD4  SORAD5  SORAD6
 75  CONTINUE
 Distribution in velocity space
   SORENI  SORENE  SORVDX  SORVDY  SORVDZ
   SORCOS  SORMAX  SORCTX  SORCTY  SORCTZ
 70 CONTINUE
\end{verbatim}

\newpage
\textbf{Meaning of the Input Variables for primary sources}

\medskip

\begin{list}{ }{}
\item[\textbf{NSTRAI }] Number of different sources (``Strata"), which
are computed
one after the other and are linearly superimposed at the end of the run.

(NSTRAI $\le$ NSTRA, see ``Parameter-Statements") \item[\textbf{INDSRC }]
\begin{list}{ }{}
\item[{INDSRC(ISTRA)=0--5}]
the input data for stratum ISTRA are read here,
but may be modified in some user routine (SAMUSR) or interface routine (INFCOP,
at entry IF2COP(ISTRA))
\item[{INDSRC(ISTRA)=6}]
no input data for stratum ISTRA are read here. The definition of this stratum
must be entirely in some problem specific routine (IF2COP, etc.).
See \cref{sec3.4} for one such example,
namely the default surface recycling
source model as specified in coupled B2-EIRENE runs.
\item[{INDSRC(ISTRA)=-1}]
the input data for stratum ISTRA are read here, and no attempt is made to
modify these. I.e.\ IF2COP(ISTRA) is not called.
\end{list}
\item[\textbf{ALLOC }]  Allocation of CPU-time to stratum weighted as

(1-ALLOC)*NPTS+ALLOC*FLUX
\item[\textbf{AMPTS}]  (available since 2014 in master version) Multiplier, used simultaneously
for permitted total CPU time NTCPU (input block 1), and for specified number of MC histories NPTS(ISTRA) (see below).
Default: AMPTS=1.0
\item[\textbf{TXTSOU }]
Text to characterize the stratum (name of the source) on the
printout file.
\item[\textbf{NLAVRP }]
= .TRUE.  not in use
\item[\textbf{NLAVRT }]
= .TRUE.  not in use
\item[\textbf{NLSYMP }]
= .TRUE.

Symmetrize profiles with respect to poloidal (y-) co-ordinate $x^2$,
i.e., with respect to the poloidal surface $x^2=$
PSURF((NP2ND+1)/2) in case NP2ND
is an odd integer, or with respect to the cell center
$x^2=$ PZONE(NP2ND/2) in case NP2ND is an even integer.
\item[\textbf{NLSYMT }]
= .TRUE.

ditto, but for toroidal (z-) co-ordinate,
i.e., for toroidal surface

TSURF((NT3RD+1)/2) or TZONE(NT3RD/2) respectively.
\item[\textbf{NPTS
}] NPTS $>$ 0: Maximum number of test particle histories.

If there is more than one stratum (NSTRAI $>$  1) then the total
CPU-time NTIME (input block 1) will be distributed proportional to
NPTS to the single strata. NPTS is the maximum number of
test-particles only if sufficient CPU time is available. Otherwise
a message ``NO FURTHER COMPUTATION TIME FOR THIS STRATUM" is
printed and the particle loop for the respective stratum is
stopped.

NPTS = 0: this stratum is ``turned off''.

NPTS $<$ 0: no limitation in the number of particles. The entire
CPU time assigned to this stratum will be used up. (NPTS is reset
to the largest integer on the machine. Hence: some care is needed
here in case of multiple strata, in combination with the
ALLOC-options to assign CPU time to individual strata).


\item[\textbf{NINITL }]
If NINITL $>$ 0 : seed for initialization of random number
generator. The results for all those individual strata can be
reproduced exactly for which the same number of test-flights is
computed as in a previous run.


If NINITL = 0 : no initialization of random numbers for the
particular stratum. In case of the first stratum, the default
initialization is used. Runs can only be reproduced, if the same
number of test-flights is computed for each stratum. Somewhat
weakened correlation between subsequent runs as compared to the
NINITL $>$ 0 option.

If NINITL $<$ 0: truly random initialization (determined by machine
clock). These runs cannot be reproduced exactly. Subsequent runs are
uncorrelated. (E.g.: recommended for stochastic approximation
procedures in nonlinear applications.
\item[\textbf{NEMODS }]
Flag to select one of the preprogrammed
source energy conditional distributions
given the source particle's position and
species.
(See: ``distribution in velocity space", below)
\item[\textbf{NAMODS }]
Flag to select one of the preprogrammed conditional
source angular distributions given the position,
species and energy of the source particle.
(See: ``distribution in velocity space", below)

\item[\textbf{NMINPTS }]
Minimum number of test particles enforced for this stratum, independent of NTCPU
flag in input block 1.  Hence: setting this flag may increase EIRENE run time above the cpu-time
assigned to a run in the first input line in input block 1.
Default: NMINPTS = 0

\item[\textbf{FLUX, SCALV}]~

\begin{list}{ }{}
\item[{SCALV=0}]
(default) FLUX = Source strength in Ampere.

FLUX is the scaling factor for all surface- or volume averaged
tallies.

FLUX is an ``atomic flux" (or: an ``atomic ion flux"). Each source
particle may carry a different flux NPRT(ISPZ) (initial weight)
depending on the species ISPZ (see below: distribution for the
species index). NPRT is specified in the blocks 4 and 5. The total
``atomic" source particle flux for each stratum is scaled to be
FLUX. For example, a $H_2$ molecule source, with NPRT$_{H_2}$ = 2,
is treated as if a flux of FLUX/1.602E-19/2 $H_2$-molecules per
second is emitted, resulting in an equivalent ``atomic flux"
FLUX/1.602E-19 per second. \item[{SCALV$\ne$0}] The default
scaling of tallies with FLUX can be overruled by this flag. The
common scaling factor for all surface- and volume averaged tallies
is determined such that one particular tally (selected by the
ISCL\dots-flags described below) has the prescribed value SCALV. This
determines the scaling of all other volume averaged and surface
averaged tallies. By this option, for example, one can set the
neutral particle density to a prescribed value in one particular
cell. Hence, one can prescribe the local Knudsen number for
nonlinear applications including neutral-neutral interactions.
 \end{list}
\item[\textbf{IVLSF}]
\begin{list}{ }{}
\item[{=1}]
The following ISCL\dots-flags select one particular volume averaged tally
\item[{=2}]
The following ISCL\dots-flags select one particular surface averaged
tally
\end{list}
\item[\textbf{ISCLS}]
species index of selected tally
\item[\textbf{ISCLT}]
tally number of selected tally (refer to \cref{tab2,tab3})
\item[\textbf{ISCL1, ISCL2, ISCL3, ISCLB, ISCLA}]~

\begin{list}{ }{}
\item[{IVLSF=1}]~

cell numbers NRCELL, NPCELL, NTCELL, NBLOCK, NACELL, respectively.

if (NPCELL = 0) or (NTCELL = 0), then ISCL1 = NCELL, the cell number in the
1-dimensional arrays (see end of \cref{sec2.2.1}).

The additional cell region is specified by NRCELL=0, NPCELL=1, NTCELL=1,
NBLOCK=NBMLT+1 (see \cref{sec2.2}) and
the proper value of NACELL.
\item[{IVLSF=2}]
to be written
\end{list}
\end{list}

\medskip

\textbf{Distribution for the species index ISP}

\medskip

\begin{list}{ }{}
\item[\textbf{NLATM }]
= .TRUE.

Atomic source. History starts in subroutine FOLNEUT with
type index ITYP=1, species index
ISP = IATM and initial weight NPRTA(IATM) (see block 4A)
\item[\textbf{NLMOL }]
= .TRUE.

Molecule source. History starts in subroutine FOLNEUT with type
index ITYP=2, species
index ISP = IMOL and initial weight NPRTM(IMOL) (see block 4B)
\item[\textbf{NLION }]
= .TRUE.

Test ion source. History starts in subroutine FOLION with type index
ITYP=3, species
index ISP = IION and initial weight NPRTI(IION) (see block 4C)
\item[\textbf{NLPHOT }]
= .TRUE.

Photon source. History starts in subroutine FOLNEUT with type index
ITYP=0, species
index ISP = IPHOT and initial weight NPRTPH(IPHOT) (see block 4D)

Not all options for direct photon sources are fully programmed. Currently
we mostly use the bulk particle volume recombination source (see next) to
simulated radiative decay from excited states as birth profile for (bound-bound)
line-photons.

\item[\textbf{NLPLS }]
= .TRUE.

Bulk ion source. Initial co-ordinates of a bulk ion with species
index ISP = IPLS and initial weight NPRTP(IPLS) (see block 5)
are generated as primary source particles,
then a surface reflection model or a volume re-combination
model is called and atoms, molecules or test ions with species index
either IATM, IMOL, IION or IPHOT
are created.
\end{list}

One and only one of these five variables must be .TRUE. .

\begin{list}{ }{}
\item[\textbf{NSPEZ }]
Species index of the source particle
\begin{list}{ }{}
\item[{1 $\le$ NSPEZ $\le$  NATMI, NMOLI, NIONI, NPLSI }]~

NSPEZ is the (fixed) species index of the source particle.
No random sampling for the species index is done.
\item[{NSPEZ $>$  NATMI, NMOLI, NIONI, NPLSI }]~

(depending upon the type of the source particle) the species index
is sampled from the distribution DATM, DMOL, DION, DPLS,
respectively.  The ``surface species distributions'' DATM, DMOL,
DION and DPLS are read in the block ``Data for General Reflection
Models", see \cref{sec2.6}.

\item[{NSPEZ $ = $ 0 }]~

the species index of the particle is directly sampled from the
``analog distribution" WEISPZ, i.e., no biased source species
sampling. WEISPZ is defined internally by the code.

The distribution WEISPZ is currently defined only for NLPLS=TRUE
sources, and here only for surface recycling sources using the STEP-function option (in SAMSRF, see further below
this section and \cref{sec2.7.1}). There it is set according to the local bulk ion flux
composition.

WEISPZ may also be transferred into a run via
user specified source sampling (SAMUSR.f, see \cref{sec3.4} )


In all other cases NSPEZ must be positive.


\item[{NSPEZ $<$  0 }]~

The ``surface species distributions'' DATM, DMOL, DION and DPLS are
considered as biased source species distributions, whereas the
analog (physical) distribution is provided automatically by the
array WEISPZ from the source sampling routines SAMPNT, SAMLNE,
SAMSRF or SAMVOL respectively. An appropriate weight correction is
carried out after sampling from DATM, DMOL, DION or DPLS,
respectively, in subroutine LOCATE.

\end{list}

Note: If a step function (function STEP, see below) is used for
sampling the start position of a test particle on a surface, then
the species index NSPEZ automatically also fixes the choice of the
index ISPZ for the spatial step function STEP(ISTEP, ISPZ, \dots)
selected by the flags SORLIM and SORIND (=ISTEP) (see below). This
default can be overruled  when SORIND has three digits.



\end{list}

\medskip

\textbf{Distribution in physical space}

\medskip
In this section options are described to sample the
starting point $\vr_0$ of a new test particle.

Together with the starting position also a vector
$\vC(\vr_0) = \vC = (CRTX,CRTY,CRTZ)$
is generated,
which may be used to distinguish one particular direction
for the
distribution in velocity space (see paragraph below:
distribution in velocity space). Internally this reference
vector $\vC$
is always normalized to unit length 1. $\vC$ is interpreted
as an outward pointing vector (at the source point).
Outflows (e.g.\ plasma fluxes onto a surface segment)
have a positive velocity component in direction of
this vector. Emitted particles (recycling, reflection,
emission from a point source) have a negative
velocity component in this reference direction $\vC$.

$\vC$ is irrelevant for isotropic velocity distribution.

It is stressed again that the reference unit vector
$\vC = (C_X,C_Y,C_Z)$ controlling un-isotropic
angular distributions
is always taken as an
\emph{outward pointing (surface-normal)
vector.} Outflowing background plasma particles
(species option NLPLS = TRUE, see above)
onto a surface element first have a positive velocity
component in direction
of $\vC$:  $\vV_b \dot \vC > 0$.
Recycled test particles then are emitted/reflected
with a negative velocity
component, i.e  $\vV_t \dot \vC < 0$.
Even in case of other sources, when no surface normal
can be assigned naturally (point sources (gas puff), volume sources)
this reference vector $\vC$ points \emph{away}
from the direction of
preferential test particle emission.

\begin{list}{ }{}
\item[\textbf{NLPNT }]
= .TRUE.     Point Source
\item[\textbf{NLLNE }]
= .TRUE.     Line Source  (not ready)
\item[\textbf{NLSRF }]
= .TRUE.     Surface Source
\item[\textbf{NLVOL }]
= .TRUE.     Volume Source
\item[\textbf{NLCNS }]
= .TRUE.     Initial conditions source (sampling from census array), for time-depen\-dent mode of operation, see input blocks 1. and 13.
\end{list}

One and only one of these 5 variables must be .TRUE.

\medskip

\textcolor[rgb]{0.00,0.00,1.00}{\textbf{Point source}
}
Flags for the ``distribution in physical space" not mentioned here
are irrelevant for point sources.
\begin{list}{ }{}
\item[\textbf{NSRFSI (=NPNTSI) }]
Total number of different points, over which the starting points
for this stratum are distributed (corresponds to ``sub-strata" option for surface and
volume sources,
there to facilitate sampling of spatial coordinates).
\end{list}
The next deck of 4 input cards is read NSRFSI times, one deck for each ``sub-stratum".
\begin{list}{ }{}
\item[\textbf{INUM }]
irrelevant; labelling index for sub-strata
\item[\textbf{SORWGT }]
Relative frequency for starting point labelled INUM.
The sum of SORWGT for all
NPNTSI points is normalized to one internally.
\item[\textbf{NRSOR }]
\begin{list}{ }{}
\item[{$>$ 0 }]
x- or radial
cell number NRCELL of the zone containing the point source.
\item[{= 0 }]
NRCELL is found automatically from the ``standard mesh" zoning.
\item[{$<$ 0 }]
only for surface sources (NLSRF), see below.
\end{list}
\item[\textbf{NPSOR }]
ditto, for y- or poloidal cell number NPCELL
\item[\textbf{NTSOR }]
ditto, for z- or toroidal cell number NTCELL
\item[\textbf{NBSOR }]
standard mesh block number NBLOCK. Defaulted to NBLOCK = 1,
if NBSOR $\le$ 0
\item[\textbf{NASOR }]
additional cell number NACELL, if point source is located
outside the standard mesh. Defaulted to NACELL = 0, if at least one
of the variables NRCELL, NPCELL or NTCELL are larger than zero.
\item[\textbf{NISOR }]
polygon index IPOLG. Meaningless if NLPLG = .FALSE.
\item[\textbf{SORAD1 }]
x-co-ordinate of source point X0
\item[\textbf{SORAD2 }]
y-co-ordinate of source point Y0
\item[\textbf{SORAD3 }]
z-co-ordinate of source point Z0
\item[\textbf{SORAD4, SORAD5, SORAD6 }] ~

are the x,y and z components of a reference directional
vector $\vC =$ (CRTX, CRTY, CRTZ)
which may be used to distinguish one particular direction for the
distribution in velocity space (see below). Internally this vector
is normalized to length 1. Irrelevant for an isotropic velocity
distribution.
\end{list}

\medskip

\textcolor[rgb]{0.00,0.00,1.00}{\textbf{Line source}}

to be written

\medskip

\textcolor[rgb]{0.00,0.00,1.00}{\textbf{Surface source}}

Flags for the ``distribution in physical space" not mentioned here are
irrelevant for surface sources.

\begin{list}{ }{}
\item[\textbf{NSRFSI }]
Total number of different surfaces, or surface segments, over which the starting points
for this stratum are distributed (``sub-strata", to facilitate sampling of spatial coordinates).
\end{list}
The next deck of 4 input cards is read NSRFSI times, one deck for each ``sub-stratum".
\begin{list}{ }{}
\item[\textbf{INUM }]
irrelevant; labelling index for sub-strata
\item[\textbf{INDIM }]
\begin{list}{ }{}
\item[{= 0 }]
source on ``additional surface" ASURF (see block 3B)
\item[{= 1 }]
source on ``standard surface" RSURF, x- (or radial) mesh

(see block 2A and 3A) \item[{= 2 }] source on ``standard surface"
PSURF, y- (or poloidal) mesh

(see block 2B and 3A) \item[{= 3 }] source on ``standard surface"
TSURF, z- (or toroidal) mesh

(see block 2C and 3A)
\item[{= 4 }]
source on a surface composed of one or more segments of radial and/or
poloidal polygons. The further details of the spatial distribution
are defined in code coupling routines, i.e., the code coupling routine
IF1COP must be
called. Special versions of IF1COP are available, e.g., in the code
segments $COUPLE_{B2}$, $COUPLE_{B2.5}$
(coupling to B2 (BRAAMS) multi-fluid plasma code), $COUPLE_{DIVIMP}$
(coupling
to DIVIMP impurity ion kinetic transport code) or  $COUPLE_{Ufile}$
(TRANSP-code format).

The position on the surface is sampled from a (piecewise constant) step function
defined from plasma fluxes
onto that surface
vs. arc-length.
The flags INSOR, INGRDA and INGRDE described below are set automatically
in this option and hence need not be specified.
\end{list}
\item[\textbf{INSOR }]
number of the surface in the mesh RSURF, PSURF, TSURF or ASURF
respectively. (Redundant in case INDIM=4)
\item[\textbf{INGRDA, INGRDE }]
same as IRPTA, IRPTE flags in input block 3a. Defines subrange
on standard surfaces, on which the source is distributed.
Irrelevant for sources on additional surfaces.
\item[\textbf{SORWGT }]
Relative frequency for starting points
on surface labelled INUM. The sum  of SORWGT for all
NSRFSI surfaces is normalized to one internally.
\item[\textbf{SORLIM }]
= KLMN

if SORLIM $\le$ 0, the user supplied
Subroutine SAMUSR
is called to sample all 3 initial
co-ordinates (X0,Y0,Z0), see \cref{sec3.4}.

if SORLIM $>$  0, then one of the preprogrammed options is used
(the digits L,M and N are relevant only for surface sources).
In this case:
\begin{list}{ }{}
\item[{N }]
Index to select one of the preprogrammed distributions
in radial or x-direction on the surface.
\item[{M }]
Index to select one of the preprogrammed distributions
in poloidal or y-direction on the surface.
\item[{L }]
Index to select one of the preprogrammed distributions
in toroidal or z-direction on the surface.
\item[{K }]
Index to select one of the preprogrammed distributions
for the starting time.
\item[{M,N,L = 0}] ~

The respective co-ordinate is computed from the 2 others and
from the equation for surface number INSOR.
\end{list}

Thus, one and only one of these 3 digits must be equal to 0,
because the
birth-point for a surface source
is determined already by two co-ordinates and the labeling                                                             o
index of the surface.

If INDIM=0, any one of the 3 digits can be the 0,
depending upon the particular equation for the surface ASURF(INSOR).

In case INDIM=1, one has to set N=0 (is now done automatically), and the poloidal (or y) and
toroidal (or z) co-ordinate is sampled according to the flags M and L.

Correspondingly in case  INDIM=2 one must
specify M=0, and  in case INDIM=3 the flag L=0 has to be set (is redundant).

\begin{list}{ }{}
\item[{L,M,N = 1 }]
$\delta$-distribution at (a+b)/2
\item[{L,M,N = 2 }]
Uniform distribution on the interval [a,b]
\item[{L,M,N = 3 }]
Truncated exponential decay
with decay length $\lambda$  on the interval [a,b].
I.e.\ the sampling distribution reads:

$f(x) = c  \cdot   \exp( - x/ \lambda)$ if $x \in [a,b]$ and $f(x) =
0$ elsewhere,

with normalizing constant

$c = \left\{ \lambda (\exp[ -a/ \lambda] - \exp[-b/
\lambda])\right\}^{-1}$
\item[{L,M,N = 4 }]
Step-function (see below: Function STEP, \cref{sec2.7.1})
(only one of either
L or M or N should be 4)
\item[{K = 1 }]
$\delta$-distribution at TIME0 for time of particle birth.  (A delta
function source in time for the kinetic equation in integral form
corresponds to an initial condition for time-dependent linear
kinetic integro-differential equation).
\item[{K = 2 }]
Uniform distribution in [TIME0,TIME0+DTIMV] for time of particle
birth.\\
Default: K=2 in time-dependent mode (NTIME $>$0) and K=1, TIME0=0 in
time-independent mode (NTIME = 0), see \cref{sec2.1}.
\end{list}
\item[\textbf{SORIND }]
Flag to choose one from the various step functions (SORLIM-option
4), which have been defined in the initialization phase. Up to NSTEP
(PARMUSR, see \cref{sec3.1}) step functions can be described
there. SORIND is the labelling index of the selected step function.

Each step function STEP(ISTEP, \dots) can consist of step functions for
fluxes of up to NSPZ species, see \cref{sec2.7.1}. NSPZ depends upon
the initialization of this function. By default the source species
index NSPEZ is used when sampling from step functions.

New option (Aug. 2006), e.g.\ for testing isotope effects: If SORIND
$\geq 100$, then the 3rd digit is used to select the species index
from step function ISTEP. I.e.: Let SORIND = LMN, then MN is used to
sample from  step function ISTEP = MN  for species ISPZ = L.
This concerns the spatial distribution. The species index itself of
the sampled particle is still determined by the flag NSPEZ, see
above.


\item[\textbf{SOREXP }]
Decay length $\lambda$  in the exponential distribution (option 3)
\item[\textbf{SORIFL }]
The first of the 4 digits can be used to overrule the default
orientation of the surface normal at the birth point, or if
ILSIDE = 0 for this particular surface.
If this digit is nonzero, a value
1 would lead to a test flight originating
from the surface as if a particle has been incident onto this surface
in the positive direction, and the value 2 means that this imaginary
particle has been striking in the negative direction.

The last 3 digits of SORIFL act as LMN of the ILSWCH flag described
in section 2.3B assuming incidence in the positive direction.

If any of this 3 digits equals zero, than the ILSWCH flag for
this particular surface is activated. (See also: section 2.3B,
input flag ILSWCH)
\item[\textbf{NRSOR }]
as for point source, but additionally:

If NRSOR  $<$   0, NRCELL is found from the step-function data (see
below: Function STEP,  \cref{sec2.7.1}) (if N is equal to 4) or
is returned from SAMUSR  (if  SORLIM  $<$   0)
\item[\textbf{NPSOR }]
as for point source, but additionally:

If NPSOR  $<$   0, NPCELL is found from the step-function data (see
below: Function STEP,  \cref{sec2.7.1}) (if M is equal to 4) or
is returned from SAMUSR  (if  SORLIM  $<$   0)
\item[\textbf{NTSOR }]
as for point source, but additionally:

If NTSOR  $<$   0, NTCELL is found from the step-function data (see
below: Function STEP,  \cref{sec2.7.1}) (if L is equal to 4) or
is returned from SAMUSR  (if  SORLIM  $<$   0)
\item[\textbf{NBSOR }]
as for point source, but additionally:

If NBSOR  $<$   0, NBLOCK is found from the step-function data (see
below: Function STEP,  \cref{sec2.7.1}) (if N is equal to 4) or
is returned from SAMUSR  (if  SORLIM  $<$   0)
\item[\textbf{NASOR }]
as for point source, but additionally:

If NASOR  $<$   0, NACELL is found from the step-function data (see
below: Function STEP, \cref{sec2.7.1}) (if N is equal to 4) or
is returned from SAMUSR  (if  SORLIM  $<$   0)
\item[\textbf{NISOR }]
as for point source, but additionally:

If NISOR $<$  0, IPOLG is found from the step-function data (see
below: Function STEP, \cref{sec2.7.1}) (if N is equal to 4) or
is returned from SAMUSR  (if  SORLIM $ < $  0)
\end{list}
The six parameters in the next card SORAD\dots are used to define the
sampling intervals in the three co-ordinate directions x, y and z.
For some geometry options and surface types (i.e., \emph{radial,
poloidal, toroidal,} or \emph{additional}) these boundaries of the
sampling interfaces, in some co-ordinates, are automatically found
from the specified range [INGRDA, INGRDE] in the corresponding grid,
and the corresponding SORAD\dots values are then not used for those
coordinates. In case of doubt see initialization phase of subroutine
SAMSRF to find out which options precisely are available, or contact
the EIRENE team at FZ-Juelich.

Only in case of \emph{additional surfaces} the information on this
deck is always fully used.
\begin{list}{ } {}
\item[\textbf{SORAD1 }]
Left endpoint a of interval [a,b] for
x- or radial co-ordinate
\item[\textbf{SORAD2 }]
Right endpoint b of interval [a,b] for
x- or radial co-ordinate
\item[\textbf{SORAD3 }]
as SORAD1, for y- or poloidal co-ordinate
\item[\textbf{SORAD4 }]
as SORAD2, for y- or poloidal co-ordinate
\item[\textbf{SORAD5 }]
as SORAD1, for z- or toroidal co-ordinate
\item[\textbf{SORAD6 }]
as SORAD2, for z- or toroidal co-ordinate
\end{list}

The reference direction unit vector $\vC = (CRTX,CRTY,CRTZ)$
for the velocity space distribution
(see paragraph below: velocity space sampling)
is, by default, the ``positive outward normal vector" of
surface INUM
at the birth point. It thus need not be specified for the
surface source option. (see ``Standard Mesh Surfaces" and/or
``Additional Surfaces").

\medskip

\textcolor[rgb]{0.00,0.00,1.00}{\textbf{Volume source }}

Flags for the ``distribution in physical space" not mentioned here are
irrelevant for volume sources.

\begin{list}{ }{}
\item[\textbf{NSRFSI }]
Total number of subregions of the standard mesh in which the starting points
for this stratum are distributed (``sub-strata", to facilitate sampling of spatial coordinates).
\end{list}
The next deck of 4 input cards is read NSRFSI times, one deck for each ``sub-stratum".
\begin{list}{ }{}
\item[\textbf{INUM }]
irrelevant; labelling index for sub-strata
\item[\textbf{INDIM, INSOR }]
not in use for volume sources
\item[\textbf{INGRDA, INGRDE }]
same as IRPTA, IRPTE flags in input block 3a. Defines subregion
of standard mesh, on which the source is distributed.
\item[\textbf{SORLIM }]
if SORLIM $\le$ 0, the user supplied
Subroutine SAMUSR
is called to sample all 3 initial
co-ordinates (X0,Y0,Z0),
see \cref{sec3.4}.

If SORLIM $>$  0, then the preprogrammed option is used.
\item[\textbf{SORIND }]
identifies volume recombination reaction IRRC for bulk plasma species IPLS,
as specified in input block 5, reaction decks for bulk ions.
(E.g., to distinguish between effects from
three-body, radiative and di-electronic recombination).

SORIND=KREC=IRRC

If more than one recombination process is specified for bulk ion
species IPLS, then IRRC is the number of one particular such
process, counted by the sequence of input in input block 5. See
printout activated by the TRCAMD flag (block 11) for the correct
value of IRRC in case of doubt.

In case SORIND=0 the sum over all relevant IRRC for the selected
background species IPLS is taken as ```recombination"' source for
IPLS. (Only for EIRENE version 2001 or younger).

Note: Both the source strength FLUX and the relative weight of subregions
(if any)
SORWGT are automatically determined from the volume source data on the
atomic data array TABRC1(IRRC,ICELL).
\item[\textbf{SORAD1 }]
unused. Parameter for user supplied sampling routine SAMUSR
\item[\textbf{SORAD2 }]
unused. Parameter for user supplied sampling routine SAMUSR
\item[\textbf{SORAD3 }]
unused. Parameter for user supplied sampling routine SAMUSR

\item[\textbf{SORAD4, SORAD5, SORAD6 }] ~

are the x,y and z components of a reference directional
vector $\vC =$ (CRTX, CRTY, CRTZ)
which may be used to distinguish one particular direction for the
distribution in velocity space (see below). Internally this vector
is normalized to length 1. Irrelevant for an isotropic velocity
distribution. (Same as for point sources, see above)
\
\end{list}

\medskip

\textbf{Distribution in velocity space}

\medskip

Together with a position (and time) $\vr_0,t_0$
of a source particle, species index $isp$,  also a specific
reference vector $\vC = (CRTX,CRTY,CRTZ)$ for sampling in
velocity space has been selected, e.g.: a surface normal
vector for surface sources, or a preferential direction of emission
from a point source or volume source, etc., see below.

Let $\vC = (C_X,C_Y,C_Z)$ be this reference directional unit vector
defined together with the point of birth of the test particle.

Furthermore, a set of local background (``plasma") data at the birth point has been provided
during spatial source sampling. I.e.:
\begin{equation}\label{local-source}
n_i(ip,\vr,t), T_e, T_i(ip,\vr,t), \vV_i(ip,\vr,t), E_i(ip,\vr,t), Sh
\end{equation}
with $ip=1, NPLS$, all background medium particle species,
for ion density, electron and ion temperature,
ion flow velocity, mean incident ion energy
and target surface sheath potential,
respectively. These local ``birth point background parameters"
are either set in case of step-function sampling, FUNCTION STEP,
or can be provided in case of user-source-sampling (SAMUSR), i.e., if  SORLIM $<$
0, or are computed from the known cell numbers and the input
background tallies at the place of birth.

Depending upon the value of the flag \textbf{NEMODS}, the following
energy distributions $f(E|\vr_0, t_0, isp)$ are available. Note that this conditional energy distribution
$f(E|\dots)$ may also
depend upon the further input parameters

\textbf{SORENI, SORENE, SORVDX, SORVDY, SORVDZ}

and is conditional on the birth position (and time) $\vr_0,t_0$
of the test particle
via local plasma parameters
$T_e ,   T_i , \vV_i, \dots $, as well as on the
selected species index $isp$ of the newly launched
test particle.

$T_i,\vV_i$ and $E_i$ may be different for different background ion
species $ip,  ip = 1,NPLS$.
The choice of a particular $ip_0=IPL$ from these
NPLS background species indices, for velocity space
sampling, if test particle $isp$ is to be launched, is described below,
digit K of the NEMODS flag.

Furthermore, let $\vV_{i,\pi}(ip_0)$ and $\vV_{i,\sigma}(ip_0)$
be the
components of the velocity drift vector $\vV_i(ip_0)$ normal
and parallel,
respectively, to a (possibly fictious) surface element
(similar  notation as in \cref{maxw-flux}), the
orientation of which is locally defined at the initial position of
the test particle by the above mentioned reference ``outward normal unit
vector" $\vC$.

In the earliest applications of EIRENE to plasma target
surface recycling sources
only the electrostatic sheath action was included.
The velocity $\vV_{i,\pi}$ corresponded then to the parallel (to a magnetic field $\vB$)
plasma flow velocity (often: sonic) onto a target surface which was assumed to be essentially
orthogonal to the magnetic field $\vB$ (e.g.\ a poloidal limiter side surface).
In case of toroidal belt limiters or poloidal divertor targets the velocity $\vV_{i,\pi}$
is interpreted as the parallel plasma flow velocity at the electrostatic sheath entrance.
I.e.\ the bending of the sonic or supersonic flow $\vV_{i,\|}$ parallel to the magnetic field
at the magnetic pre-sheath entrance
to a sonic flow $\vV_{i,\pi}$ normal to the target entering the electrostatic sheath,
by the action of the magnetic pre-sheath field, is not carried out inside EIRENE.
Check \cref{sec1.5} for generalizations, e.g.\ to also include magnetic pre-sheath models inside EIRENE.



\begin{list}{ }{}
\item[\textbf{NEMODS }]
= KLMN

The choice of one of the following energy distributions is
made depending upon the value of the first digit N of NEMODS:
\begin{list}{ }{}
\item[{N = 1 }]
Mono-energetic source $f(E) = \delta (E-E0)$ with fixed energy

E0 = SORENI (\si{\electronvolt})
\item[{N = 2 }]
Mono-energetic source $f(E) = \delta (E-E0)$ with energy

E0 = SORENI $\cdot T_i(ip_0) $ + SORENE $\cdot T_e $

\item[{N = 3 }]
(only for surface sources)
As N=2, but with added sheath acceleration, see \cref{sec2.7.2} below


\item[{N = 4 }]
Mono-energetic source f(E) = $\delta$ (E-E0) with energy

E0 = EMAX($T_i(ip_0) ,V_{i,\pi}(ip_0) ,V_{i,\sigma}(ip_0)$)

where EMAX is the mean energy from a truncated (one-sided) shifted Maxwellian flux
at temperature $T_i(ip_0)$ drifting by a velocity

$(V_{i,\pi}(ip_0), V_{i,\sigma}(ip_0)$).

One obtains (see \cref{eq1.30}, omitting unnecessary indices) in \cref{sec1}) by integration EMAX = $\gamma_E   T_i(ip_0) $

with $\gamma_E   =   (
{\cal V}_{\pi}^2  +   2   +  {\cal V}_{\sigma}^ 2   )   +   0.5
\frac{g({\cal V}_{\pi} ) - 1}{g ({\cal V}_{\pi} )}$ ,

${\cal V}   = V / \sqrt{ \frac{{2 T_i}}{{m_i}} }$, and $g(x)   =   1 +
\sqrt{\pi} x  [ 1 + \mbox{erf}(x) ]   \exp(x^2 )$.

(The normalized velocities ${\cal V}, {\cal V}_{\pi}, {\cal V}_{\sigma}$ used here coincide with the
isothermal Mach number if $T_i(ip_0) = T_e$, see again \cref{maxw-flux}).
\item[{N = 5 }]
(only for surface sources)
As N=4, but with added sheath acceleration, see \cref{sec2.7.2} below

E0 = EMAX($T_i(ip_0) ,V_{i,\pi}(ip_0) ,V_{i,\sigma}(ip_0)$) +
ESHEAT
\item[{N = 6 }]

\emph{surface source}

The velocity vector $\vV_0 = (VX_0,VY_0,VZ_0)$ is sampled from a
truncated Maxwellian flux $f_{flux}= \vV \cdot \vC f_{max} (T_i(ip_0) ,\vV_i(ip_0) , \vC )$
with a drifting Maxellian $f_{max}$ at temperature
$ T_i(ip_0) $ and drift velocity $\vV_i(ip_0) $. The (one-sided) flux is
across a surface element described locally by the outward normal
vector $\vC $. The mean energy from this sampling distribution is
that given by option N=4.

\emph{point or volume source}

The velocity vector $\vV_0 = (VX_0,VY_0,VZ_0)$ is sampled from a
truncated Maxwellian density distribution $f_{max} (T_i(ip_0) ,\vV_i(ip_0) , \vC
)$ at temperature $ T_i(ip_0) $ and which is shifted by a velocity $\vV_i(ip_0) $

\item[{N = 7 }]
(only for surface sources)
As N=6, but with added sheath acceleration, see \cref{sec2.7.2} below

\item[{N = 8 }]
Mono-energetic source $f(E) = \delta (E-E0)$ with energy

$E0 = E_i(ip_0)$, with $E_i(ip_0)$ being the local mean incident ion energy, for example defined in connection with the spatial
``step-function" option, see \cref{sec2.7.1}.
\item[{N = 9 }]
(only for surface sources)
As N=8, but with added sheath acceleration, see \cref{sec2.7.2} below
\end{list}
The choice of temperature parameters $T_e$ and $T_i(ip_0)$
in the energy distributions N = 4, 5, 6 or 7
of the plasma fluids entering the sheath region
is controlled by the second digit M of NEMODS:
\begin{list}{ }{}
\item[{M =0 }]
$T_e$ is the local electron temperature taken from input tally TEIN
on the grid.

In case of NLPLS, the default parameter $T_i(ip_0)$ is the local ion
temperature for bulk ion species $ip_0=IPL$, IPL is the species index ISP of
the source particle itself, by default.

In case of NLION, the default parameter $T_i(ip_0)$ is the local ion
temperature for bulk ion species $ip_0=IPL$, IPL is determined such that
the mass and charge number and the charge state of the source
particle IION match the mass, charge number and charge state of the
bulk ion IPL. If no such bulk ion is found, $T_i(ip_0) = 0$.

These default choices of background medium species index $ip_0=IPL$ can be overruled by the K digit of this
input flag, see below.

In case of .not.(NLPLS.or.NLION), $T_i(ip_0)$ =0.
\item[{M =1 }]
The local ion temperature
$T_i$ is replaced by the input constant SORENI (\si{\electronvolt}) and
$T_e$ is replaced by the constant SORENE (\si{\electronvolt}) in options
N = 4 to N = 7.
\item[{M =2 }]
not in use
\item[{M =3 }]
$T_e$ is the local electron temperature, and $T_i$ is the local
background ion temperature of ion species $ip_0 = IPL=K$, K is the fourth digit
of the NEMODS flag, see below.

\end{list}
The choice of velocity parameters $\vV$ in the energy distributions
N=4,5,6 or 7 is controlled by the third digit L of NEMODS:
\begin{list}{ }{}
\item[{L =0 }]
In case of NLPLS, the parameter $\vV$ is the local ion drift velocity
for bulk ion species IPL, IPL is the species index of the source particle.

In case of .NOT.NLPLS, $\vV$ =0.
\item[{L =1 }]
The local drift velocity vector
$\vV $ is replaced by the constant vector (SORVDX, SORVDY,
SORVDZ) (\si{\centi\metre\per\second}), regardless of the point of birth.

\item[{L =2 }]
The local drift velocity vector
$\vV $ is replaced by the constant vector
(SORVDX, SORVDY, SORVDZ) $ \cdot $ CS, where CS is the isothermal ion sound
velocity for the species and at the temperatures
chosen above (\si{\centi\metre\per\second}).

Therefore, { (SORVDX,SORVDY,SORVDZ) }
are now understood as Mach numbers in x,y and z direction
respectively.
\item[{L =3 }]
$\vV$ is the local
background ion drift velocity of ion bulk species IPL=K, K is the fourth digit
of the NEMODS flag, see below.

\item[{K options }] ~

The digit K (unless zero) is used to specify one of the
background ion species IPL, from which the $T_i(IPL), \vV(IPL) $
parameters in the sampling distributions described above are taken.

\end{list}
\item[\textbf{SORCOS, SORMAX }]~

Depending upon the value of the
flag \textbf{NAMODS} (see below) various different angular distributions
may be selected. Each one depends upon the two parameters P = SORCOS and
Q = SORMAX.
\item[\textbf{SORCTX, SORCTY, SORCTZ }]~

If this input vector is NOT zero:

The unit outward pointing reference vector $\vC $ mentioned
above is overwritten and explicitly now given by normalization of
SORCTX, SORCTY, SORCTZ.

Else (default):

As described above this vector $\vC$ is internally found as the
outward surface normal vector (in case of surface sources)
at $\vr_0$ or (in case of point, line, or volume sources)
by the input parameters $\vC =(SORAD4,SORAD5,SORAD6)$.

By this option, for example, the mean direction of emission
from
a surface or a gas puff can be influenced.
In case of surface sources the distribution
of the polar angle may then not necessarily be centered
around the inner
surface normal vector any longer, but may be tilted.

\begin{list}{ }{}
\item[{NAMODS = 1 }]~

The polar angle $\vartheta$ against the reference unit vector
$\vC = (C_X,C_Y,C_Z)$ of the source particle's velocity is sampled from
a cosine**P distribution around the ``inner normal vector" (-1.0)
$\cdot \vC $, i.e., $f( \vartheta )   d \vartheta   \sim \sin(
\vartheta )  \cdot  \cos^p ( \vartheta )   d \vartheta$.

Important special cases:
\begin{list}{ }{}
\item[{P = 0 }]
isotropic distribution
\item[{P = 1 }]
cosine distribution
\item[{P $\gg$ 1 }]
close to $\delta$-distribution around vector $-1  \cdot \vC $
\end{list}
\item[{NAMODS = 2 }]~

The polar angle against $-1  \cdot  \vC $
is sampled from a
Gaussian distribution with zero mean value, and
the parameter P now is used
for the standard deviation (degree) of that distribution.
\end{list}
The second parameter Q is the cut-off angle (degree)
for the polar angle distribution
\begin{list}{ }{}
\item[{note:  }]
Q $\le 180^{\circ}$  is enforced internally

Q $\le 90^{\circ}$ is enforced internally for surface sources

Q = 0  for a beam, i.e., for an angular $\delta$-distribution at
$-1  \cdot  \vC $
\end{list}


\end{list}
In all pre-programmed
cases the azimuthal angle (around the axis $\vC$)
is equally distributed from
$0^{\circ}$ to  $360^{\circ}$.

\subsection{Piecewise constant ``Step-functions" for sampling}
\label{sec2.7.1}
The statement:

\texttt{A = STEP(NSPZI,NSPZE,NSMAX(ISTEP),ISTEP)}

initializes (piecewise constant) step functions for random sampling by the inversion method,
ISTEP is an index labeling one such a step function, and there may be NSTEP different step functions (species index)
with one common set of abscissae.

The statement:

\texttt{B = STEP1(IINDEX,ISTEP,RNF,ISPEZ)}

converts a uniformly distributed random number RNF into a random number sampled
from the step function  no. ISTEP.

\medskip

Sampling of the birth-point of a trajectory is sometimes done from such
piecewise constant functions (\texttt{EIRENE} function \texttt{STEP}(\dots) ), e.g.,
if one digit of the flag \texttt{SORLIM} has been set to the value 4.
This is the case if, for example, a recycling ion flux distribution is given
in discretized form on a grid along a target surface (from an external \ac{CFD} plasma solver).

The
random sampling of a coordinate $u$ is from the step-function
\texttt{FLSTEP(RRSTEP)}. There is storage for up to \texttt{NSTEP} (see \texttt{PARMUSR})
such step-functions.

$u_i = RRSTEP(i,ISTEP), i = 1, \dots, NSMAX(ISTEP) - 1$, is a discrete ordered (either increasing or decreasing) set of
abscissae $u_i$, at which the piecewise constant step-function has
a jump to $FLSTEP(u_i)$.

The final interval ends at \texttt{RRSTEP(NSMAX(ISTEP),ISTEP)}

Internally the piecewise constant function $FLSTEP$ is translated into a (normalized) cummulative
distribution $VF(u)$ for random number generation. FLSTEP itself is not necessarily normalized.


Let RNF be a uniform random number on the interval
\begin{equation}\nonumber
0.0 =\texttt{VF(RRSTEP(1,ISTEP))} \rightarrow \texttt{VF(RRSTEP(NSMAX),ISTEP)} = 1.0,
\end{equation}

or on a sub-interval thereof.

By calls to the entry \texttt{STEP1(IINDEX,ISTEP,RNF,ISPEZ)} of Function \texttt{STEP}
the randomly sampled index $i$=\texttt{IINDEX} is returned, as well as a coordinate $u$ sampled from an
uniform distribution on the $i^{th}$ increment $\Delta u_i$
= \texttt{RRSTEP(i+1, \dots) - RRSTEP(i, \dots)}: $u$=\texttt{STEP1}

In case of \texttt{LEVGEO=1} or \texttt{LEVGEO=2} $u$ has the simple meaning of one of the
spatial coordinates in an \texttt{EIRENE} run.

However, $u$ need not necessarily be one of the 3 spatial
coordinates. \texttt{RRSTEP} can also stand, e.g., for an arc-length along
a polygon (\texttt{LEVGEO=3}), or for the cumulated length of an element in
a selected list of sides of triangles (\texttt{LEVGEO=4}) or for the
cumulated area of a surface element in a selected list of surfaces
of tetrahedrons (\texttt{LEVGEO=5}). In these latter cases only the index
\texttt{IINDEX} is used, and a point from an uniform distribution on the
line-, or  surface element, to which \texttt{IINDEX} points, must
still be sampled.


The probability density functions \texttt{FLSTEP} for sampling a position
(or cell index) \texttt{RR}
are set in function \texttt{STEP} from a piecewise constant distribution
function (e.g., from a surface flux density distribution). In the same call to \texttt{STEP} also their corresponding
cumulative functions \texttt{VF} are obtained by integration of \texttt{FLSTEP} over \texttt{RRSTEP}, and normalization.

These step-functions have to be set in the initialization phase
(by calls to function \texttt{STEP} from \texttt{SAMUSR}, \texttt{INFUSR}, etc.).


The initialization call to function \texttt{STEP} reads:

\texttt{FLUX = STEP(NSPZI,NSPZE,NSMAX(ISTEP),ISTEP)}

This call will initialize the step-function no. \texttt{ISTEP}, for a range of species
\texttt{ISPZ=NSPZI, NSPZE}.

At this point the grid RRSTEP and the distribution density FLSTEP on that grid must be defined (module CSTEP).
RRSTEP is an ordered set of numbers, the abscissae of the points when FLSTEP jumps from one value to another.
RRSTEP may have dimension length, area, volume or any other dimension of the independent variable for the
piecewise constant distribution FLSTEP.

In the initialization call to function STEP the distribution density \texttt{FLSTEP} (Flux per unit length or area or volume,
or whatever the dimension of RRSTEP is.)
is then converted into
a (set of) cumulative distribution $VF$ by
\begin{equation}\nonumber
VF(ISPZ,ISTEP,i)= \sum_{j=1}^{j=i} \Delta u_j \cdot FLSTEP(ISPZ,ISTEP,j)
\end{equation}
with, again, $\Delta u_J=RRSTEP(J+1,ISTEP)-RRSTEP(J,ISTEP)$,

one for each species index separately. This cumulative flux-distribution is then in units of a flux (source strength), typically in EIRENE: in [\si{\ampere}].

After this the total flux (normalization constant) \texttt{FLTOT(ISPZ,ISTEP)} = \texttt{VF(ISPZ,ISTEP,NSMAX)} is stored,
and then the cumulative distributions $VF$ are normalized
to
\begin{equation}\nonumber
VF(\dots,i) = VF(\dots,i)/FLTOT(\dots),  \ \ \ \ i=1,NSMAX(\dots).
\end{equation}
hence enforcing $VF(ISPZ,ISTEP,NSMAX) = 1$ for random sampling.


The pdf FLSTEP, the cumulative distributions \texttt{VF}
and normalization factors \texttt{FLTOT} themselves
are stored in module \texttt{CSTEP.f},
as well as the normalized invers functions
(quantile functions) of the cumulative distributions,
which are needed for efficient random sampling.



Default initializations of functions \texttt{STEP} are carried out in the
initialization of surface source sampling routine \texttt{SAMSRF}, entry
\texttt{SMSRF0}.

Random sampling using such step-functions \texttt{STEP} is done, e.g., from
surface source sampling routine \texttt{SAMSRF}, entry \texttt{SMSRF1}.

In addition to the value of the distribution
\texttt{FLSTEP(i)}, the normalization constants \texttt{FLTOT} and the normalized cumulative \texttt{VF(i)} in each interval
some further quantities may be defined as function of the interval index
``$i$":

\begin{list}{ }{}
\item[{TESTEP }]
electron temperature at the surface segment ``$i$"
\item[{TISTEP }]
ion temperature at the surface segment ``$i$" (array of length NPLS )
\item[{DISTEP }]
ion density at the surface segment ``$i$" (array of length NPLS)
\item[{IRSTEP }]
cell number in 1st (x or radial) grid of the surface segment ``$i$".
\item[{IPSTEP }]
cell number in 2nd (y or poloidal) grid of the surface segment ``$i$".
\item[{ITSTEP }]
cell number in 3rd (z or toroidal) grid of the surface segment ``$i$".
\item[{IBSTEP }]
block number, see section 2e
\item[{IASTEP }]
additional cell number of the surface segment ``$i$"
\item[{ELSTEP }]
ion energy flux at the surface segment ``$i$" (array of length NPLS)
\item[{SHSTEP }]
electrostatic sheath potential at the  surface segment ``$i$"
\end{list}

\textbf{Default step functions}

Step functions for sampling the radial coordinate (at a given or sampled  poloidal and toroidal position)
(\texttt{INDIM}=1, \texttt{SORLIM}-digit N = 4) are defined internally during the initialisation step of surface sampling routine \texttt{SAMSRF} if the selected step functions
\texttt{SORIND=ISTEP} has not been defined externally already.
Also in case of unstructured grids (\texttt{LEVGEO}=4,5) default step functions can be defined for non-default surfaces.


In such cases the particle flux  spatial sampling distribution \texttt{FLSTEP} of background species $ip$ on a radial grid segment $i$ or triangle side $i$ or tetrahedron side $i$ is computed as
\begin{equation}
FLSTEP(ip,i)=n_i(ip) \times cs(ip)
\end{equation}
with $n_i(ip)$ the local ion density of species $ip$ (=DISTEP(ip,i)) and $cs(ip)$ is the isothermal acoustic speed of ion species
$ip$, i.e.
\begin{equation}
cs(ip)=\sqrt{[Z_i T_e+T_i(ip)]/m_i(ip)}.
\end{equation}
Here $Z_i$ denotes the ion electric charge.

The mean incident ion energy $ELSTEP=E_i$ corresponds to that resulting from the \texttt{NEMOD} = 2, 3 options,
with $E_i=3.0 T_i + 0.5 Z_i T_e$ and the sheath potential $SHSTEP= Sh$ is set to the sheath potential flag for the non-default surface: \texttt{FSHEAT}, see input block 6.

Note that by setting SHSTEP and hence parameter $Sh$, the evaluation of the full sheath potential function $\Delta \Phi$
in energy sampling options \texttt{NEMODS}, digit N=3,5,7 and 9 is overruled by using $\Delta \Phi = Sh$ instead.


\subsection{Electrostatic sheath acceleration}
\label{sec2.7.2}

For surface sources, and ionic source particles (test ions or bulk ion, with charge $Z_i$) an acceleration of the
sampled velocity vector $\vv_0$ towards the surface, in the direction normal to the target surface, may be added.
I.e.\ the velocity space sampling described above (\texttt{NEMODS} options, digit N) is regarded to provide the velocity at the entrance
of the electrostatic sheath in front of a target surface.

Energy sampling options \texttt{NEMODS} N= 3,5,7,9 add this sheath contribution to velocities sampled from the corresponding options N= 2,4,6, and 8, respectively.

Let $\vv_0=\vV_0$. This sheath acceleration is achieved by setting a new velocity $\vV_1$:

${\vV}_1   =
{\vV}_0   + {\vV}_S$, ${\vV}_S   =
V_S
   \cdot   \vC $ and $V_S$ evaluated such that
\begin{equation}
m_i    \cdot    V_S \left
( {\vV}_0
\vC \right )   +   \frac{m_i}{2}
V_S^2   =
ESHEAT(n_i,{\vec{v}}_{\pi,i} ,T_e),
\end{equation}
i.e., the energy $E$ of the particle is increased from $E_0$ to $E_1$ by the amount ESHEAT (\si{\electronvolt}) as
compared to the option without sheath acceleration, and this increased energy comes from an additional speed component $V_s$ normal towards the target surface.

The sheath acceleration $ESHEAT = Z_i \cdot \Delta \Phi \cdot T_e$  is computed from the sheath voltage
$\Delta \Phi$, where $Z_i$ is
the electric charge of the sampled particle at birth point, $Z_i=0$ for neutral particles.
The sheath potential drop $\Delta \Phi$ is
given by \cref{sheath1} in \cref{sec1}, unless overruled by parameter $Sh>0.$, (see \cref{local-source}, in which latter case $ESHEAT=Z_i \cdot Sh \cdot T_e$.

If $Sh\leq 0$ and the sheath potential according to equation \cref{sheath1} cannot be found either (e.g.\ because $\vV_{\pi,i}(ip) \equiv 0.)$, then
$Sh = FSHEAT(MSURF)$ (the input sheath parameter for source surface no. MSURF, input blocks 3 and 6) is used.
